% !TEX program = xelatex
% =====================================================================
%  BÁO CÁO ĐỒ ÁN TRÍ TUỆ NHÂN TẠO, HỌC PHẦN AIP202
%  Khoa Công nghệ Thông tin, Trường Đại học Kiến trúc Đà Nẵng
%
%  Build:  latexmk -xelatex do-an.tex
%  Overleaf: Menu > Compiler > XeLaTeX, rồi bấm Recompile.
%
%  Khung tài liệu chuyển thể từ mẫu đồ án của Trường Đại học Nha Trang
%  (github.com/nd-hung/thesis-template, tác giả TS. Nguyễn Đình Hưng).
%  Toàn bộ nhận diện, phông chữ, bảng màu, trang bìa và cấu trúc chương
%  được thay theo chuẩn của DAU và theo đề cương học phần AIP202.
% =====================================================================
% \documentclass phải nằm ở file này, không được chuyển vào settings.tex,
% vì Overleaf dò file chính bằng cách tìm file có chứa \documentclass.
\documentclass[12pt, a4paper, oneside]{book}
% !TEX root = ../do-an.tex
% =====================================================================
%  styles/settings.tex
%  Các thiết lập chung, nạp từ file chính do-an.tex.
%  SINH VIÊN KHÔNG SỬA FILE NÀY.
%
%  Lưu ý: lệnh khai báo lớp tài liệu nằm ở file do-an.tex chứ không nằm
%  ở đây. Lý do: Overleaf dò file chính bằng cách tìm chuỗi khai báo lớp
%  tài liệu trong các file. Nếu chuỗi đó xuất hiện ở file này, Overleaf
%  sẽ biên dịch nhầm file này và báo lỗi "no legal end found", vì file
%  này chỉ chứa thiết lập chứ không chứa phần thân tài liệu.
% =====================================================================

% Toàn bộ nhận diện DAU: phông, màu, bìa, hộp trình bày, header
\usepackage{styles/dau}
\dauLayout

% Giãn dòng và khoảng cách đoạn theo quy cách đồ án
\linespread{1.25}
\usepackage[skip=8pt plus 1pt, indent=15pt]{parskip}

% Ghi chú: mẫu gốc dùng gói atbegshi để vứt trang đầu, vì trang bìa của
% mẫu gốc dựng bằng minipage nên sinh ra một trang trắng. Trang bìa ở
% đây dựng bằng môi trường titlepage nên không có trang trắng, và nếu
% giữ thủ thuật đó thì chính trang bìa sẽ bị vứt mất.


% =====================================================================
%  PHẦN EM ĐIỀN: thông tin đồ án
% =====================================================================
% Loại báo cáo
\crname{BÁO CÁO ĐỒ ÁN TRÍ TUỆ NHÂN TẠO}
% Học phần
\cmonhoc{Học phần AIP202 \textperiodcentered{} 2 tín chỉ}
% Tên đề tài. Dùng \\ để ngắt dòng cho cân đối trên bìa.
\ctname{XÂY DỰNG HỆ THỐNG HỖ TRỢ GỢI Ý\\ĐƯỜNG ĐI CỦA CƠN BÃO VIỆT NAM}
% Họ tên sinh viên và mã số sinh viên
% Nhóm hai sinh viên: bìa mặc định in "tên (MSSV)" cho một người, nên ở đây
% in bảng tên | MSSV rồi nuốt phần " (MSSV)" mà mẫu tự thêm phía sau.
\def\boquaMSSV#1(#2){}
\cstuname{\begin{tabular}[t]{@{}l@{\hspace{0.15cm}}l@{}}
  Nguyễn Viết Mạnh & 2351220241\\
  Nguyễn Đình Phú  & 2351220247
\end{tabular}\boquaMSSV}
\cmssv{2351220241, 2351220247}
% Lớp
\clopsv{23CT5}
% Ngành
\csCouncil{Công nghệ thông tin}
% Giảng viên hướng dẫn
\csSupervise{TS. Đỗ Phúc Hảo}
% Năm học
\cttime{2025--2026}
% Khoa quản lý
\csdeptname{KHOA CÔNG NGHỆ THÔNG TIN}
% Nơi thực hiện
\csplace{Đà Nẵng}

% Gói parskip chèn khoảng cách đoạn vào ô cột X, làm lệch dòng tiêu đề bảng
\AtBeginEnvironment{tabularx}{\setlength{\parskip}{0pt}}

\begin{document}

% Bìa trang trọng có khung kép. Muốn dùng bìa hiện đại theo bộ giáo
% trình của khoa thì thay dòng dưới bằng \coverpagehiendai
\coverpage

% ---------------------------------------------------------------------
%  PHẦN ĐẦU: đánh số trang bằng số La Mã thường
% ---------------------------------------------------------------------
\frontmatter
\dauPageStyle
% Chân trang một dòng cho nhóm hai sinh viên
\newcommand{\chantrangSV}{\fancyfoot[L]{\footnotesize\color{gray500}Nguyễn Viết Mạnh \textperiodcentered{} Nguyễn Đình Phú}}
\chantrangSV

% ---------------------------------------------------------------------
%  LỜI CAM ĐOAN
% ---------------------------------------------------------------------
\begin{declaration}

Chúng tôi xin cam đoan đồ án này là công trình do chính chúng tôi thực hiện
dưới sự hướng dẫn của \dauGVHD. Các nội dung nghiên cứu, số liệu và kết quả
trình bày trong đồ án là trung thực, được sinh ra từ mã nguồn của dự án và
chưa từng được công bố trong bất kỳ công trình nào khác.

Mọi số liệu, hình ảnh, mã nguồn và kết quả của tác giả khác được sử dụng
trong đồ án đều được trích dẫn nguồn gốc rõ ràng trong phần tài liệu tham
khảo. Dữ liệu quỹ đạo bão lấy từ NOAA IBTrACS và dữ liệu tái phân tích lấy từ
NCEP/NCAR, đều là dữ liệu công khai. Toàn bộ phần mềm và thư viện được dùng
đều có giấy phép cho phép sử dụng. Những phần có sử dụng công cụ trí tuệ nhân
tạo hỗ trợ đã được nêu rõ trong báo cáo.

Nếu phát hiện có bất kỳ sự gian lận nào, chúng tôi xin chịu hoàn toàn trách
nhiệm trước Khoa và Nhà trường.

\vspace{1.4em}
\noindent\okyten{Nguyễn Viết Mạnh}{Sinh viên thực hiện}\hfill
\okyten{Nguyễn Đình Phú}{Sinh viên thực hiện}

\end{declaration}

% ---------------------------------------------------------------------
%  LỜI CẢM ƠN
% ---------------------------------------------------------------------
\begin{acknowledgments}

Chúng em xin gửi lời cảm ơn chân thành đến \dauGVHD, người đã định hướng đề
tài, góp ý và theo sát tiến độ của chúng em trong suốt quá trình thực hiện đồ
án. Chúng em cũng xin cảm ơn quý thầy cô Khoa Công nghệ Thông tin, Trường Đại
học Kiến trúc Đà Nẵng đã trang bị nền tảng kiến thức để chúng em hoàn thành
công việc này. Chúng em cảm ơn Cơ quan Khí quyển và Đại dương Quốc gia Hoa Kỳ
(NOAA) đã công bố miễn phí bộ dữ liệu IBTrACS và tái phân tích NCEP/NCAR, là
nguồn dữ liệu chính của đồ án. Cuối cùng, chúng em xin cảm ơn gia đình và bạn
bè đã luôn động viên chúng em trong thời gian học tập.

\vspace{1.4em}
\noindent\okyten{Nguyễn Viết Mạnh}{Sinh viên thực hiện}\hfill
\okyten{Nguyễn Đình Phú}{Sinh viên thực hiện}

\end{acknowledgments}

% ---------------------------------------------------------------------
%  TÓM TẮT
% ---------------------------------------------------------------------
\begin{abstract}

Việt Nam chịu ảnh hưởng trung bình nhiều cơn bão và áp thấp nhiệt đới mỗi
năm, nên biết sớm hướng đi của bão giúp chủ động phòng tránh. Đồ án xây dựng
một hệ thống trí tuệ nhân tạo gợi ý đường đi của cơn bão, phát biểu dưới dạng
bài toán dự báo chuỗi thời gian nhiều bước: từ bốn quan trắc gần nhất (24
giờ), hệ thống ước lượng vị trí tâm bão sau 6, 12, 24, 48 và 72 giờ kèm nón
bất định. Đồ án huấn luyện và so sánh bảy phương pháp, gồm hai mốc so sánh
(persistence, CLIPER), k láng giềng gần nhất dựa trên các cơn bão lịch sử
tương tự (k-NN analog), gradient boosting, LSTM, Attention-LSTM và
Transformer, trên 1.089 cơn bão ở Tây Bắc Thái Bình Dương giai đoạn
1980--2024, chia tập theo năm để không rò rỉ thông tin. Trên tập kiểm thử
2021--2024, các phương pháp học giảm sai số 72 giờ khoảng 200~km so với
persistence (676~km xuống 462--476~km), có ý nghĩa thống kê. Thí nghiệm cho
thấy học độ dịch chuyển thay vì tọa độ tuyệt đối giảm sai số 6 giờ từ 115~km
xuống 33~km, và bổ sung đặc trưng dòng dẫn đường khí quyển cho mô hình kết hợp
LSTM và Transformer giảm thêm 35~km tại 72 giờ. Hệ thống được triển khai thành
API và bản đồ web, mỗi gợi ý kèm danh sách cơn bão lịch sử tương tự để người
dùng hiểu vì sao có gợi ý đó. Kết quả cho thấy với dữ liệu chỉ có quỹ đạo, mô
hình học sâu không vượt trội phương pháp cổ điển; nút thắt nằm ở thông tin khí
quyển chứ không ở kiến trúc mạng.

\vspace{0.8em}
\noindent\textbf{Từ khoá:} dự báo quỹ đạo bão, chuỗi thời gian, LSTM,
Transformer, k láng giềng gần nhất, lập luận dựa trên tình huống.

\end{abstract}


\tableofcontents
\listoffigures
\listoftables
% Bỏ dòng dưới nếu đồ án của em không có giải thuật nào
\listofalgorithms

% ---------------------------------------------------------------------
%  THÂN BÁO CÁO: năm chương (04 và 05 cùng thuộc Chương 4)
% ---------------------------------------------------------------------
\mainmatter
\dauPageStyle
\chantrangSV

% ---------------------------------------------------------------------
%  CHƯƠNG 1 — MỞ ĐẦU
% ---------------------------------------------------------------------
\chapter{Mở đầu}
\label{ch:modau}

\begin{muctieu}
Chương này nêu bối cảnh và lý do chọn đề tài, phát biểu bài toán, các mục
tiêu có thể kiểm chứng, phạm vi thực hiện, đóng góp chính và bố cục của báo
cáo.
\end{muctieu}

\section{Bối cảnh và lý do chọn đề tài}

Việt Nam có hơn 3.200~km bờ biển nằm cạnh Biển Đông, một trong những vùng
xoáy thuận nhiệt đới hoạt động mạnh nhất thế giới. Phần lớn các cơn bão ảnh
hưởng đến nước ta hình thành ở Tây Bắc Thái Bình Dương, đi vào Biển Đông rồi
đổ bộ vào đất liền; bão YAGI năm 2024 đi thẳng vào miền Bắc là một ví dụ gần
đây. Với công tác phòng chống thiên tai, câu hỏi quan trọng nhất là tâm bão
sẽ ở đâu trong một đến ba ngày tới, vì câu trả lời quyết định nơi nào cần sơ
tán và khi nào phải cấm biển.

Dự báo nghiệp vụ hiện nay dựa trên mô hình số trị khí quyển, cần siêu máy
tính và đội ngũ chuyên môn. Trong khi đó, hơn bốn mươi năm dữ liệu quỹ đạo
bão đã được công bố công khai và có chất lượng kiểm soát tốt. Từ đó đề tài
đặt ra hai câu hỏi:
\begin{itemize}[leftmargin=1.6em, itemsep=2pt]
  \item Một hệ thống trí tuệ nhân tạo chỉ học từ dữ liệu lịch sử có thể gợi ý
        đường đi của bão chính xác đến mức nào?
  \item Các nhóm phương pháp AI khác nhau, từ lập luận dựa trên tình huống
        đến học máy và học sâu, khác nhau ra sao khi đặt trên cùng một bài
        toán, cùng một tập dữ liệu và cùng một cách đánh giá?
\end{itemize}

Đề tài được chọn vì ba lý do. Thứ nhất, bài toán có ý nghĩa thực tiễn rõ
ràng với Việt Nam. Thứ hai, nó cho phép vận dụng nhiều nhóm kiến thức của
học phần trên một bộ dữ liệu thật. Thứ ba, kết quả đo được bằng một đại
lượng dễ hiểu là số kilômét lệch so với vị trí thực, nên có thể so sánh
khách quan giữa các phương pháp.

Có thể hình dung độ khó của bài toán qua phương pháp đơn giản nhất là
\term{persistence}, tức giả định bão giữ nguyên hướng và tốc độ của 6 giờ gần
nhất. Trên dữ liệu kiểm thử của đồ án, cách ngoại suy này chỉ lệch khoảng
31~km sau 6 giờ, nhưng lệch tới 676~km sau 72 giờ, dài hơn cả quãng đường
chim bay từ Đà Nẵng ra Hà Nội. Nói cách khác, ở hạn ngắn bão gần như đi theo
quán tính, còn ở hạn dài thì quán tính không còn đủ. Khoảng trống giữa hai con
số đó là chỗ để một phương pháp học từ dữ liệu chứng minh giá trị, và cũng là
thước đo mà đồ án dùng để tự đánh giá.

Một điểm cần nói rõ ngay từ đầu là cách hiểu chữ ``gợi ý'' trong tên đề tài.
Cơn bão không lựa chọn đường đi, nên đây không phải bài toán tìm đường trên đồ
thị. Hệ thống ước lượng vị trí tương lai của tâm bão, nói rõ mức độ không chắc
chắn của ước lượng đó và chỉ ra căn cứ lịch sử của nó. Người đọc kết quả vẫn
là người ra quyết định; hệ thống chỉ cung cấp thêm một nguồn tham khảo.

\section{Các hướng tiếp cận hiện có và khoảng trống}

Dự báo quỹ đạo bão hiện có ba hướng chính. Hướng thứ nhất là mô hình số trị:
giải hệ phương trình động lực học khí quyển trên lưới ba chiều. Đây là hướng
chính xác nhất và là nền tảng của dự báo nghiệp vụ, nhưng đòi hỏi siêu máy
tính, dữ liệu đồng hóa toàn cầu và đội ngũ chuyên gia, vượt xa điều kiện của
một đồ án học phần. Hướng thứ hai là các phương pháp thống kê cổ điển, tiêu
biểu là CLIPER~\cite{Neumann1972}, kết hợp quán tính với khí hậu học; chúng rẻ
và ổn định nên thường được dùng làm mốc so sánh. Hướng thứ ba là học máy trên
dữ liệu best-track như IBTrACS~\cite{Knapp2010}, từ các phương pháp tìm bão
tương tự đến mạng hồi quy~\cite{Hochreiter1997} và
Transformer~\cite{Vaswani2017}.

Khi khảo sát hướng thứ ba, nhóm thấy ba khoảng trống mà một đồ án học phần có
thể đóng góp:
\begin{itemize}[leftmargin=1.6em, itemsep=2pt]
  \item Nhiều kết quả chỉ báo cáo một mô hình học sâu so với persistence, không
        đặt cạnh các phương pháp đơn giản hơn trên cùng dữ liệu. Vì vậy khó biết
        phần cải thiện đến từ kiến trúc hay chỉ từ việc có học.
  \item Cách chia dữ liệu ngẫu nhiên theo quan trắc làm rò rỉ thông tin giữa
        huấn luyện và kiểm thử, vì các quan trắc liên tiếp của một cơn bão rất
        giống nhau. Điểm số thu được khi đó đẹp hơn thực tế.
  \item Kết quả thường dừng ở một con số sai số, thiếu vùng bất định được hiệu
        chỉnh và thiếu lời giải thích để người không chuyên có thể tin hoặc nghi
        ngờ một gợi ý cụ thể.
\end{itemize}
Ba khoảng trống này định hướng cho các mục tiêu và cách đánh giá của đồ án.

\section{Phát biểu bài toán}

Cho chuỗi quan trắc gần nhất của một cơn bão, gồm vị trí tâm, gió và khí áp
đo cách nhau 6 giờ, hệ thống cần gợi ý vị trí tâm bão tại các mốc 6, 12, 24,
48 và 72 giờ tiếp theo. Mỗi gợi ý đi kèm một vùng bất định cho biết mức độ
tin cậy, và danh sách các cơn bão lịch sử tương tự để người dùng hiểu gợi ý
dựa trên căn cứ nào. Chương~\ref{ch:thietke} phát biểu chi tiết đầu vào, đầu
ra và tiêu chí chạy đúng.

Về mặt học máy, đây là bài toán hồi quy chuỗi thời gian nhiều bước: từ một cửa
sổ quan trắc có độ dài cố định, mô hình xuất đồng thời năm cặp vĩ độ và kinh
độ. Chất lượng được đo bằng khoảng cách trên mặt cầu giữa vị trí gợi ý và vị
trí thực, tính bằng kilômét, và bằng tỉ lệ vị trí thực rơi vào trong vùng bất
định.

\section{Câu hỏi nghiên cứu và giả thuyết}

Từ bối cảnh và các khoảng trống nêu trên, đồ án đặt ra ba câu hỏi nghiên cứu.
Mỗi câu hỏi đi kèm một giả thuyết ban đầu và một cách kiểm chứng cụ thể
(Bảng~\ref{tab:cauhoi}). Việc ghi giả thuyết trước khi chạy thí nghiệm giúp
phần đánh giá ở Chương~\ref{ch:ketqua} không bị uốn theo kết quả.

\begin{table}[H]
\centering
\caption{Câu hỏi nghiên cứu, giả thuyết ban đầu và cách kiểm chứng.}
\label{tab:cauhoi}
\begin{tabularx}{\textwidth}{>{\bfseries\color{daunavy}}l X X X}
\toprule
\rowcolor{gray100}\bfseries\color{daunavy}Mã & \leavevmode\bfseries\color{daunavy}Câu hỏi & \leavevmode\bfseries\color{daunavy}Giả thuyết ban đầu & \leavevmode\bfseries\color{daunavy}Cách kiểm chứng \\
\midrule
CH1 & Học từ dữ liệu lịch sử có tốt hơn ngoại suy quán tính không, và từ hạn dự báo nào?
    & Có, và lợi ích tăng theo hạn dự báo.
    & So sai số từng mốc với persistence; bootstrap theo cơn bão tại 72 giờ. \\
CH2 & Mô hình học sâu có vượt các phương pháp đơn giản hơn không?
    & Có, vì mạng hồi quy và Transformer mô hình hóa chuỗi tốt hơn.
    & So sánh bảy phương pháp trên cùng cửa sổ kiểm thử, ba seed, có khoảng tin cậy. \\
CH3 & Muốn giảm sai số thêm thì nên thay đổi mô hình hay thay đổi thông tin đầu vào?
    & Chưa rõ; cần thí nghiệm có đối chứng.
    & Thử lần lượt từng thay đổi, chọn trên validation, chạy kiểm thử một lần. \\
\bottomrule
\end{tabularx}
\end{table}

Kết quả cho thấy giả thuyết của CH1 đúng, giả thuyết của CH2 sai, và CH3 có
câu trả lời rõ ràng là thông tin đầu vào. Việc một giả thuyết bị bác bỏ được
giữ nguyên trong báo cáo, vì đó là phát hiện có giá trị nhất của đồ án.

\section{Mục tiêu của đồ án}

Đồ án đặt ra năm mục tiêu, mỗi mục tiêu đều có tiêu chí để kiểm chứng ở
Chương~\ref{ch:ketqua}:
\begin{itemize}[leftmargin=1.6em, itemsep=2pt]
  \item \textbf{MT1.} Xây dựng quy trình dữ liệu từ NOAA IBTrACS: làm sạch,
        tạo đặc trưng, chia tập theo năm không rò rỉ thông tin.
  \item \textbf{MT2.} Huấn luyện và so sánh ít nhất bốn phương pháp AI với hai
        mốc so sánh trên cùng tập kiểm thử, có kiểm định thống kê.
  \item \textbf{MT3.} Mô hình tốt nhất giảm sai số vị trí tại 72 giờ ít nhất
        20\% so với phương pháp persistence.
  \item \textbf{MT4.} Cung cấp vùng bất định (nón) có độ phủ thực tế xấp xỉ
        mức danh nghĩa 80\%.
  \item \textbf{MT5.} Triển khai mô hình thành API và giao diện bản đồ, mỗi gợi
        ý kèm lời giải thích bằng các cơn bão lịch sử tương tự.
\end{itemize}

Các mục tiêu được viết sao cho có thể trả lời ``đạt'' hoặc ``không đạt'' bằng
số liệu. Bảng~\ref{tab:tieuchi} nêu tiêu chí kiểm chứng của từng mục tiêu và
nơi trình bày bằng chứng trong báo cáo.

\begin{table}[H]
\centering
\caption{Tiêu chí kiểm chứng của từng mục tiêu.}
\label{tab:tieuchi}
\begin{tabularx}{\textwidth}{>{\bfseries\color{daunavy}}l X >{\raggedright\arraybackslash}p{3.6cm}}
\toprule
\rowcolor{gray100}\bfseries\color{daunavy}Mục tiêu & \leavevmode\bfseries\color{daunavy}Tiêu chí kiểm chứng & \leavevmode\bfseries\color{daunavy}Bằng chứng \\
\midrule
MT1 & Không cơn bão nào xuất hiện ở hai tập; bộ chuẩn hóa chỉ fit trên tập huấn luyện; mỗi quan trắc truy được về tệp nguồn. & Mục thiết kế dữ liệu, Chương~\ref{ch:thietke} \\
MT2 & Mọi phương pháp được đánh giá trên cùng các cửa sổ; mạng nơ-ron chạy ba seed; có khoảng tin cậy 95\%. & Bảng~\ref{tab:ketqua-all}, \ref{tab:bootstrap} \\
MT3 & Sai số 72 giờ của mô hình tốt nhất không quá 80\% sai số của persistence. & Bảng~\ref{tab:ketqua-all}, \ref{tab:steering} \\
MT4 & Tỉ lệ vị trí thực nằm trong nón 80\% trên tập kiểm thử nằm trong khoảng 75--90\%. & Mục~\ref{sec:non} \\
MT5 & Gọi được API để nhận quỹ đạo, nón và bão tương tự; xem được trên bản đồ; bộ kiểm thử tự động chạy đạt. & Chương~\ref{ch:caidat}, mục kiểm thử phần mềm \\
\bottomrule
\end{tabularx}
\end{table}

\section{Phạm vi và đối tượng}

Đối tượng nghiên cứu là quỹ đạo tâm của xoáy thuận nhiệt đới ở Tây Bắc Thái
Bình Dương, giới hạn trong vùng 0--30°N, 100--140°E, giai đoạn 1980--2024.
Đồ án dự báo vị trí tâm bão tại năm mốc 6, 12, 24, 48 và 72 giờ.

Đồ án \textbf{không} dự báo cường độ, mưa, lũ hay nước dâng; không phát hành
cảnh báo chính thức và không gửi thông báo tới người dùng. Hệ thống chỉ phục
vụ học tập và nghiên cứu, không thay thế bản tin của Trung tâm Dự báo Khí
tượng Thủy văn Quốc gia (NCHMF).

Vùng 0--30°N, 100--140°E được chọn vì nó bao trọn Biển Đông và phần phía tây
của Tây Bắc Thái Bình Dương, nơi các cơn bão đi qua trước khi ảnh hưởng đến
Việt Nam. Giai đoạn bắt đầu từ năm 1980 vì từ đó vệ tinh địa tĩnh đã phủ
thường xuyên khu vực, nên vị trí tâm bão trong dữ liệu đáng tin cậy hơn các
thập kỷ trước. Năm mốc dự báo trải từ rất ngắn (6 giờ) đến ba ngày (72 giờ),
là khoảng thời gian thường được dùng để chuẩn bị ứng phó.

Trong phạm vi đó, đồ án dành riêng một nhóm đánh giá cho \textbf{bão ảnh hưởng
Việt Nam}, định nghĩa là cơn bão có quỹ đạo đi vào trong 300~km tính từ đường
bờ biển. Nhóm này có 27 cơn bão trong giai đoạn kiểm thử và là nhóm mà đề tài
quan tâm nhất.

\section{Phương pháp thực hiện}

Đồ án được thực hiện theo sáu bước nối tiếp, mỗi bước cho ra một sản phẩm
kiểm tra được trước khi sang bước sau:
\begin{enumerate}[leftmargin=1.8em, itemsep=2pt]
  \item \textbf{Thu thập và làm sạch dữ liệu.} Tải IBTrACS, chuẩn hóa đơn vị,
        khử trùng lặp, lọc theo vùng và ghi lại nguồn gốc từng quan trắc.
  \item \textbf{Xây dựng đặc trưng và chia tập.} Viết một bộ tạo đặc trưng dùng
        chung cho huấn luyện và suy luận; chia tập theo năm, giữ nguyên cả cơn
        bão.
  \item \textbf{Dựng mốc so sánh.} Cài đặt persistence và CLIPER-proxy trước
        khi huấn luyện bất kỳ mô hình nào, để có ngưỡng mà mọi phương pháp học
        phải vượt qua.
  \item \textbf{Huấn luyện và so sánh.} Huấn luyện k-NN analog, gradient
        boosting và ba mạng nơ-ron; chọn cấu hình trên tập validation; chạy tập
        kiểm thử một lần; kiểm định bằng bootstrap theo cơn bão.
  \item \textbf{Phân tích và thí nghiệm bổ sung.} Thí nghiệm loại bỏ để kiểm
        chứng lựa chọn thiết kế; thử các hướng giảm sai số; nghiên cứu từng
        cơn bão cụ thể.
  \item \textbf{Đóng gói và triển khai.} Xuất gói mô hình có mã băm, dựng API
        và bản đồ web, viết bộ kiểm thử tự động.
\end{enumerate}

Hai nguyên tắc được giữ xuyên suốt. Thứ nhất, mọi con số trong báo cáo đều
sinh ra từ mã nguồn trong kho và có lệnh tái lập ở Chương~\ref{ch:caidat}.
Thứ hai, tập kiểm thử không tham gia vào bất kỳ quyết định nào; mọi lựa chọn
đều dựa trên tập validation.

\section{Đối tượng sử dụng}

Hệ thống hướng tới ba nhóm người dùng, với nhu cầu khác nhau:
\begin{itemize}[leftmargin=1.6em, itemsep=2pt]
  \item \textbf{Người học và giảng viên} dùng hệ thống như một ví dụ trọn vẹn
        về quy trình học máy trên dữ liệu thật: từ dữ liệu thô đến mô hình,
        đánh giá và triển khai.
  \item \textbf{Người quan tâm đến thời tiết} xem quỹ đạo gợi ý, vùng bất định
        và các cơn bão lịch sử tương tự trên bản đồ, kèm lời nhắc rằng bản tin
        chính thức thuộc về NCHMF.
  \item \textbf{Người phát triển} gọi API để lấy kết quả dạng có cấu trúc, mỗi
        phản hồi kèm phiên bản mô hình và nguồn dữ liệu.
\end{itemize}

\section{Đóng góp chính}

\begin{itemize}[leftmargin=1.6em, itemsep=2pt]
  \item Một quy trình dữ liệu tái lập được, có truy vết nguồn gốc và chia tập
        theo năm để tránh rò rỉ.
  \item Một phép so sánh công bằng giữa bảy phương pháp thuộc bốn nhóm (mốc
        so sánh, AI cổ điển, học máy, học sâu), có khoảng tin cậy bootstrap
        theo cơn bão.
  \item Bằng chứng thực nghiệm rằng nút thắt của bài toán nằm ở thông tin đầu
        vào chứ không ở kiến trúc mô hình: bổ sung dòng dẫn đường khí quyển
        giúp giảm sai số có ý nghĩa thống kê, còn đổi kiến trúc thì không.
  \item Một hệ thống chạy được từ đầu đến cuối gồm gói mô hình có kiểm tra
        toàn vẹn, API và bản đồ web, kèm bộ kiểm thử tự động.
\end{itemize}

\section{Ý nghĩa của đề tài}

\textbf{Về học thuật}, đồ án đặt bốn nhóm phương pháp của học phần lên cùng
một bài toán và cùng một cách đánh giá: lập luận dựa trên tình huống, học máy
trên cây quyết định, mạng hồi quy và cơ chế attention. Kết quả cho thấy một
điều hay bị bỏ qua khi học về học sâu: với dữ liệu vài nghìn mẫu và đầu vào
ngắn, phương pháp đơn giản có thể ngang bằng mô hình phức tạp, và chất lượng
thông tin đầu vào quan trọng hơn kiến trúc.

\textbf{Về thực tiễn}, hệ thống cho thấy có thể dựng một công cụ tham khảo về
đường đi của bão chỉ từ dữ liệu công khai và máy tính cá nhân; lần huấn luyện
LSTM chuẩn chỉ mất khoảng hơn một phút trên CPU. Công cụ này không thay thế
dự báo nghiệp vụ, nhưng phù hợp cho giảng dạy và cho việc thử nhanh các ý
tưởng mới.

\textbf{Về phương pháp}, đồ án áp dụng một số thói quen của nghiên cứu nghiêm
túc ở quy mô nhỏ: chia tập chống rò rỉ, chọn mô hình trên validation, báo cáo
khoảng tin cậy, và nêu rõ cả những kết quả không như mong đợi.

\section{Bố cục báo cáo}

Phần còn lại của báo cáo gồm năm chương và một phụ lục.

Chương~\ref{ch:coso} trình bày cơ sở lý thuyết: cách đo sai số trên mặt cầu,
các mốc so sánh, bốn nhóm phương pháp được dùng, cách ước lượng vùng bất định
và các công cụ phần mềm.

Chương~\ref{ch:thietke} phân tích yêu cầu và trình bày thiết kế: kiến trúc
tổng thể, thiết kế dữ liệu và đặc trưng, thiết kế các mô hình, thiết kế phần
phục vụ, giao diện và quy trình chọn mô hình.

Chương~\ref{ch:caidat} mô tả việc cài đặt: tổ chức mã nguồn, các đoạn mã cốt
lõi, lệnh tái lập thí nghiệm và bản demo.

Chương~\ref{ch:ketqua} trình bày kết quả: sai số của từng phương pháp, kiểm
định thống kê, độ phủ của nón, các thí nghiệm bổ sung, nghiên cứu trường hợp,
kiểm thử phần mềm, bàn luận và hạn chế.

Chương~\ref{ch:ketluan} tổng kết những gì đã làm được, bài học rút ra và
hướng phát triển.

% ---------------------------------------------------------------------
%  CHƯƠNG 2 — CƠ SỞ LÝ THUYẾT VÀ CÔNG NGHỆ
% ---------------------------------------------------------------------
\chapter{Cơ sở lý thuyết và công nghệ}
\label{ch:coso}

\begin{muctieu}
Chương này trình bày các kiến thức nền tảng gồm bài toán dự báo chuỗi thời
gian, các phương pháp trí tuệ nhân tạo được dùng, cách đo sai số, cùng danh
mục công cụ và thư viện để xây dựng hệ thống.
\end{muctieu}

\section{Kiến thức nền tảng}

\subsection{Dự báo quỹ đạo bão như bài toán chuỗi thời gian}

\begin{dinhnghia}[Dự báo quỹ đạo nhiều mốc]
Cho chuỗi quan trắc $\mathbf{x}_{t-T+1}, \dots, \mathbf{x}_t$ của một cơn bão,
mỗi $\mathbf{x}_i \in \mathbb{R}^F$ là vec-tơ đặc trưng tại thời điểm $i$.
Bài toán là học hàm $f$ ánh xạ chuỗi này thành vị trí tâm bão
$\hat{\mathbf{p}}_{t+h} = (\hat\varphi_{t+h}, \hat\lambda_{t+h})$ tại các mốc
$h \in \mathcal{H} = \{6, 12, 24, 48, 72\}$ giờ, với $\varphi$ là vĩ độ và
$\lambda$ là kinh độ.
\end{dinhnghia}

Cơn bão không ``chọn'' đường đi như tác tử trong bài toán tìm đường; nó bị
dòng khí quyển quy mô lớn (dòng dẫn đường, \en{steering flow}) đẩy đi. Vì vậy
``gợi ý đường đi'' ở đây là bài toán \term{hồi quy nhiều đầu ra}: ước lượng vị
trí tương lai và độ bất định của nó.

\subsection{Sai số khoảng cách vòng lớn}

Sai số dự báo được đo bằng khoảng cách trên mặt cầu Trái Đất (công thức
Haversine) giữa vị trí dự báo và vị trí thực:
\begin{equation}
d = 2R \arcsin\sqrt{\sin^2\frac{\Delta\varphi}{2}
    + \cos\varphi_1\cos\varphi_2\sin^2\frac{\Delta\lambda}{2}},
\label{eq:haversine}
\end{equation}
với $R = 6371$~km là bán kính Trái Đất. Đây là độ đo chuẩn trong đánh giá dự
báo quỹ đạo bão, có đơn vị km nên dễ hiểu với người dùng.

\subsection{Các phương pháp mốc so sánh}

\term{Persistence} giả định bão giữ nguyên vận tốc của 6 giờ gần nhất.
\term{CLIPER} \en{(CLImatology and PERsistence)}~\cite{Neumann1972} kết hợp quán
tính với khí hậu học; đồ án dùng một phiên bản rút gọn (CLIPER-proxy) với
persistence có suy giảm. Một phương pháp AI chỉ có giá trị nếu thắng được hai
mốc này.

\subsection{Lập luận dựa trên tình huống: k láng giềng gần nhất}

Phương pháp \term{analog} là dạng lập luận dựa trên tình huống
\en{(case-based reasoning)}~\cite{RussellNorvig2021}: tìm trong thư viện các
tình huống bão lịch sử giống tình huống hiện tại nhất rồi dùng diễn biến tiếp
theo của chúng để dự báo. Khoảng cách giữa hai tình huống $q$ và $c$ là khoảng
cách Euclid trên các nhóm đặc trưng $g$ (đoạn quỹ đạo, vị trí, mùa) đã chuẩn
hóa, mỗi nhóm có trọng số $w_g$ và chia cho căn bậc hai số chiều $n_g$ của nhóm
để các nhóm so sánh được với nhau:
\begin{equation}
D(q, c) = \sqrt{\sum_{g} \frac{w_g^2}{n_g} \,\bigl\lVert \mathbf{z}_g(q) - \mathbf{z}_g(c) \bigr\rVert_2^2} .
\label{eq:knn-dist}
\end{equation}
Dự báo là trung bình có trọng số nghịch đảo khoảng cách của $k$ láng giềng:
$\hat{\mathbf{y}} = \sum_{j=1}^{k} \alpha_j \mathbf{y}_j$, với
$\alpha_j \propto 1 / (D_j + \epsilon)$. Ưu điểm lớn của k-NN là giải thích
được: hệ thống nêu tên các cơn bão lịch sử đã dùng.

\subsection{Gradient boosting}

Gradient boosting~\cite{Friedman2001} xây dựng tuần tự nhiều cây quyết định,
mỗi cây học phần dư của tổ hợp các cây trước. Đồ án dùng bộ hồi quy
\en{histogram gradient boosting} của scikit-learn (300 vòng, tốc độ học 0,05,
tối đa 31 lá mỗi cây), mỗi mốc và mỗi tọa độ một mô hình, đầu vào là cửa sổ
đặc trưng đã trải phẳng.

\subsection{Mạng LSTM}

LSTM \en{(Long Short-Term Memory)}~\cite{Hochreiter1997} là mạng nơ-ron hồi quy
có các cổng quên $f_t$, cổng vào $i_t$ và cổng ra $o_t$ điều khiển ô nhớ $c_t$:
\begin{equation}
\begin{aligned}
c_t &= f_t \odot c_{t-1} + i_t \odot \tanh(W_c[\mathbf{h}_{t-1}, \mathbf{x}_t] + b_c),\\
\mathbf{h}_t &= o_t \odot \tanh(c_t),
\end{aligned}
\end{equation}
nhờ đó giữ được thông tin qua nhiều bước. Trạng thái ẩn cuối cùng
$\mathbf{h}_t$ được đưa qua tầng tuyến tính để sinh ra dự báo cho cả năm mốc
cùng lúc. Biến thể \term{Attention-LSTM}~\cite{Bahdanau2015} thay $\mathbf{h}_t$
bằng tổng có trọng số của mọi trạng thái ẩn, trọng số do cơ chế attention học.

\subsection{Transformer}

Transformer~\cite{Vaswani2017} bỏ cơ chế hồi quy, dùng self-attention để mỗi
bước thời gian ``nhìn'' trực tiếp mọi bước khác:
\begin{equation}
\mathrm{Attention}(Q, K, V) = \mathrm{softmax}\!\left(\frac{QK^\top}{\sqrt{d_k}}\right) V,
\end{equation}
cộng thêm mã hóa vị trí để giữ thứ tự thời gian. Đồ án dùng phần encoder.

\subsection{Huấn luyện và kiểm định}

Mạng nơ-ron được tối ưu bằng Adam~\cite{Kingma2015} với hàm mất mát Smooth~L1,
ít nhạy với ngoại lai hơn bình phương sai số. Để so sánh hai phương pháp,
đồ án dùng \term{bootstrap}~\cite{Efron1993} lấy mẫu lại theo cơn bão, vì các
dự báo của cùng một cơn bão tương quan với nhau; khoảng tin cậy 95\% không
chứa 0 nghĩa là khác biệt có ý nghĩa thống kê.

\begin{ghinho}
Với bài toán dự báo, phải chia tập theo thời gian và giữ nguyên cả cơn bão.
Nếu chia ngẫu nhiên theo điểm quan trắc, mô hình sẽ ``thấy'' các điểm khác của
cùng cơn bão khi huấn luyện, và sai số báo cáo sẽ thấp hơn thực tế.
\end{ghinho}

\section{Công cụ và thư viện}

\begin{table}[H]
\centering
\caption{Công cụ và thư viện sử dụng trong đồ án.}
\label{tab:congcu}
\begin{tabularx}{\textwidth}{>{\bfseries\color{daunavy}}l >{\color{daumaroon}}l X}
\toprule
\rowcolor{gray100}\bfseries\color{daunavy}Thành phần & \bfseries\color{daumaroon}Phiên bản & \leavevmode\bfseries\color{daunavy}Vai trò \\
\midrule
Python        & 3.11 & Ngôn ngữ lập trình chính \\
PyTorch       & 2.x  & Huấn luyện LSTM, Attention-LSTM, Transformer \\
scikit-learn  & 1.x  & Gradient boosting, k-NN, chuẩn hóa \\
pandas, NumPy & 2.x  & Xử lý dữ liệu quỹ đạo dạng bảng \\
xarray        & ---  & Đọc dữ liệu tái phân tích NCEP/NCAR \\
FastAPI       & ---  & API phục vụ dự báo \\
Leaflet       & ---  & Bản đồ web hiển thị quỹ đạo và nón \\
pytest        & ---  & Kiểm thử tự động \\
\bottomrule
\end{tabularx}
\end{table}

\section{Các công trình liên quan}

Dự báo quỹ đạo bão truyền thống đi theo ba hướng. Hướng thống kê với CLIPER
\cite{Neumann1972} là mốc so sánh chuẩn suốt nhiều thập kỷ. Hướng mô hình số
trị khí quyển cho độ chính xác cao nhất nhưng cần siêu máy tính. Hướng học máy
gần đây dùng mạng hồi quy và Transformer học trực tiếp từ dữ liệu best-track
như IBTrACS~\cite{Knapp2010}, đôi khi bổ sung trường khí quyển từ tái phân tích
\cite{Kalnay1996}. Đồ án kế thừa hướng học máy, nhưng khác ở chỗ so sánh công
bằng nhiều họ phương pháp trên cùng tập kiểm thử có kiểm định thống kê, tập
trung vào nhóm bão ảnh hưởng Việt Nam, và dùng k-NN analog để giải thích gợi ý.

% ---------------------------------------------------------------------
%  CHƯƠNG 3 — PHÂN TÍCH VÀ THIẾT KẾ
% ---------------------------------------------------------------------
\chapter{Phân tích và thiết kế}
\label{ch:thietke}

\begin{muctieu}
Chương này phân tích yêu cầu của bài toán, sau đó lần lượt trình bày thiết kế
kiến trúc hệ thống, thiết kế dữ liệu và đặc trưng, thiết kế các mô hình học,
thiết kế phần phục vụ và quy trình chọn mô hình.
\end{muctieu}

\section{Phân tích bài toán và yêu cầu}

\subsection{Đầu vào và đầu ra}

\textbf{Đầu vào} là 4 quan trắc gần nhất của một cơn bão, cách nhau 6 giờ,
tức 24 giờ lịch sử. Mỗi quan trắc gồm thời điểm UTC, vĩ độ, kinh độ, gió, khí
áp và cấp cường độ.

\textbf{Đầu ra} gồm ba phần: vị trí tâm bão tại +6, +12, +24, +48 và +72 giờ;
bán kính nón bất định tại từng mốc; và danh sách 5 cơn bão lịch sử tương tự
nhất, dùng làm lời giải thích cho gợi ý.

\subsection{Yêu cầu chức năng}

\begin{itemize}[leftmargin=1.6em, itemsep=2pt]
  \item Nạp và làm sạch dữ liệu best-track, tạo đặc trưng, huấn luyện mô hình.
  \item Dự báo quỹ đạo qua API cho bão đang hoạt động và cho chuỗi quan trắc
        do người dùng gửi lên.
  \item Hiển thị quỹ đạo dự báo, nón bất định và bão tương tự trên bản đồ.
\end{itemize}

\subsection{Yêu cầu phi chức năng}

\begin{itemize}[leftmargin=1.6em, itemsep=2pt]
  \item \textbf{Không rò rỉ dữ liệu} giữa huấn luyện và kiểm thử.
  \item \textbf{Tái lập được}: cố định seed, lưu cấu hình và mã băm dữ liệu
        của từng lần chạy.
  \item \textbf{Toàn vẹn}: mô hình được đóng gói kèm mã băm, phần phục vụ từ
        chối gói bị sửa đổi hoặc sai phiên bản.
  \item \textbf{Minh bạch}: mọi kết quả hiển thị kèm phiên bản mô hình, nguồn
        dữ liệu và cảnh báo ``chỉ mang tính tham khảo''.
\end{itemize}

\subsection{Tiêu chí chạy đúng}

Hệ thống được coi là đạt khi sai số Haversine trên tập kiểm thử thấp hơn
persistence có ý nghĩa thống kê, và nón 80\% phủ khoảng 80\% vị trí thực.

\section{Thiết kế kiến trúc hệ thống}

Kiến trúc tổng thể được trình bày ở Hình~\ref{fig:kientruc}. Hệ thống tách
thành hai nửa có vòng đời khác nhau. Nửa huấn luyện chạy ngoại tuyến, đi từ
dữ liệu thô đến gói mô hình. Nửa phục vụ chạy trực tuyến, nạp gói mô hình đó
để trả lời yêu cầu dự báo. Hai nửa chỉ giao tiếp qua gói mô hình, nên có thể
huấn luyện lại mà không sửa phần phục vụ.

\begin{figure}[H]
\centering
\resizebox{\textwidth}{!}{%
\begin{tikzpicture}[node distance=10mm and 9mm]
  \node[objbox] (data) {IBTrACS WP\\ 1980--2024};
  \node[clsbox, right=of data] (clean) {Làm sạch,\\ chuẩn hóa};
  \node[clsbox, right=of clean] (feat) {Tạo đặc trưng\\ (47 + 11 biến)};
  \node[clsbox, right=of feat] (split) {Cửa sổ chuỗi,\\ chia theo năm};
  \node[clsbox, right=of split] (train) {Huấn luyện,\\ chọn trên val};
  \node[objbox, right=of train] (bundle) {Gói mô hình\\ + thư viện analog};
  \node[objbox, below=of feat] (env) {NCEP/NCAR R1\\ (dòng dẫn đường)};
  \node[clsbox, below=of bundle] (api) {FastAPI\\ dự báo + nón};
  \node[objbox, left=of api] (web) {Bản đồ web\\ + bão tương tự};
  \draw[flowarr] (data) -- (clean);
  \draw[flowarr] (clean) -- (feat);
  \draw[flowarr] (feat) -- (split);
  \draw[flowarr] (split) -- (train);
  \draw[flowarr] (train) -- (bundle);
  \draw[flowarr] (env) -- (feat);
  \draw[flowarr] (bundle) -- (api);
  \draw[flowarr] (api) -- (web);
\end{tikzpicture}}
\caption{Sơ đồ kiến trúc tổng thể của hệ thống.}
\label{fig:kientruc}
\end{figure}

Mã nguồn chia thành các mô-đun độc lập, mỗi mô-đun ứng với một khối trong sơ
đồ:
\begin{itemize}[leftmargin=1.6em, itemsep=2pt]
  \item \code{ingestion}: tải và phân tích dữ liệu nguồn, ghi lại nguồn gốc.
  \item \code{features}: làm sạch, tạo đặc trưng, chuẩn hóa.
  \item \code{datasets}: cắt cửa sổ chuỗi và chia tập.
  \item \code{models}, \code{training}: định nghĩa mô hình, vòng huấn luyện,
        checkpoint và đóng gói.
  \item \code{evaluation}: tính sai số và bootstrap.
  \item \code{inference}: nạp gói mô hình, k-NN analog, nón bất định.
  \item \code{api}: giao diện HTTP và phục vụ trang bản đồ.
\end{itemize}

\section{Thiết kế dữ liệu}

\subsection{Nguồn và làm sạch}

Nguồn chính là NOAA IBTrACS v04r01~\cite{Knapp2010}, lưu vực Tây Bắc Thái Bình
Dương (WP), giai đoạn 1980--2024. Quy trình làm sạch gồm bốn bước: chuẩn hóa
đơn vị gió từ knot sang m/s; khử bản ghi trùng lặp; chỉ giữ quan trắc synop 6
giờ (00, 06, 12, 18 UTC); và giới hạn vùng 0--30°N, 100--140°E. Sau làm sạch
còn \textbf{24.024 quan trắc của 1.089 cơn bão}. Mỗi quan trắc lưu nguồn, URL
và mã băm SHA-256 của tệp gốc, nên có thể truy ngược bất kỳ con số nào về tệp
dữ liệu đã sinh ra nó.

\subsection{Chia tập không rò rỉ}

Các quan trắc liên tiếp của cùng một cơn bão rất giống nhau. Nếu chia ngẫu
nhiên theo quan trắc, mô hình sẽ gặp lại gần như cùng một tình huống ở cả
huấn luyện lẫn kiểm thử và cho điểm số đẹp giả tạo. Vì vậy dữ liệu được chia
\textbf{theo năm, giữ nguyên cả cơn bão} (Bảng~\ref{tab:chiatap}). Bộ chuẩn
hóa đặc trưng chỉ được fit trên tập huấn luyện rồi dùng lại nguyên vẹn cho
validation, test và khi suy luận.

\begin{table}[H]
\centering
\caption{Chia tập dữ liệu theo năm.}
\label{tab:chiatap}
\begin{tabularx}{\textwidth}{>{\bfseries\color{daunavy}}l >{\color{daumaroon}}c c X}
\toprule
\rowcolor{gray100}\bfseries\color{daunavy}Tập & \bfseries\color{daumaroon}Giai đoạn & \bfseries\color{daunavy}Số bão & \leavevmode\bfseries\color{daunavy}Mục đích \\
\midrule
Huấn luyện  & 1980--2016 & 895 & Học tham số ($\approx$7.800 cửa sổ) \\
Validation  & 2017--2020 & 110 & Chọn mô hình, siêu tham số, bán kính nón \\
Kiểm thử    & 2021--2024 & 84  & Đánh giá một lần (714 cửa sổ) \\
Ảnh hưởng VN & 2021--2024 & 27 & Tập con của kiểm thử: đi vào trong 300~km bờ biển (278 cửa sổ) \\
\bottomrule
\end{tabularx}
\end{table}

\subsection{Đặc trưng}

Một bộ tạo đặc trưng duy nhất (\code{FeatureBuilder}) dùng chung cho huấn
luyện và suy luận, nhờ đó hai giai đoạn không thể lệch nhau về cách tính. Mỗi
quan trắc có 47 đặc trưng, tóm tắt ở Bảng~\ref{tab:dactrung}. Thí nghiệm mở
rộng ở mục~\ref{sec:steering} bổ sung 11 đặc trưng dòng dẫn đường lấy từ tái
phân tích NCEP/NCAR R1~\cite{Kalnay1996}.

\begin{table}[H]
\centering
\caption{Các nhóm đặc trưng đầu vào.}
\label{tab:dactrung}
\begin{tabularx}{\textwidth}{>{\bfseries\color{daunavy}}l X}
\toprule
\rowcolor{gray100}\bfseries\color{daunavy}Nhóm & \leavevmode\bfseries\color{daunavy}Đặc trưng \\
\midrule
Vị trí     & Vĩ độ, kinh độ, chênh lệch và độ trễ 1--3 bước \\
Động học   & Quãng đường bước, hướng di chuyển, tốc độ (km/h), gia tốc, góc rẽ, tốc độ rẽ \\
Thời gian  & sin/cos của giờ, ngày, tháng, ngày trong năm; cờ mùa bão \\
Địa lý     & Khoảng cách, hướng tới bờ biển Việt Nam và các điểm tham chiếu (Hoàng Sa, Trường Sa, Đà Nẵng, \dots) \\
Cường độ   & Gió (m/s), khí áp (hPa), biến thiên của chúng, mã one-hot cấp bão \\
Dòng dẫn đường (mở rộng) & Dòng dẫn đường lớp sâu trong vành 200--800~km, gió 500~hPa, đứt gió 200--850~hPa, độ cao 500~hPa và chênh lệch 5° về bốn hướng \\
\bottomrule
\end{tabularx}
\end{table}

\section{Thiết kế mô hình}

\subsection{Biểu diễn đầu ra: độ dịch chuyển}

Thay vì để mạng dự báo trực tiếp tọa độ tuyệt đối, mô hình dự báo
\term{độ dịch chuyển} so với vị trí quan trắc cuối $\mathbf{p}_t$:
\begin{equation}
\hat{\mathbf{p}}_{t+h} = \mathbf{p}_t + \Delta\hat{\mathbf{p}}_h .
\end{equation}
Vị trí hiện tại $\mathbf{p}_t$ là thông tin đã biết chính xác. Nếu học tọa độ
tuyệt đối trên đầu vào đã chuẩn hóa, mạng phải tốn sức ``khôi phục'' lại chính
vị trí này. Học độ dịch chuyển cho phép mạng tập trung vào phần thực sự cần
dự báo là hướng và tốc độ di chuyển. Mục~\ref{sec:ablation} kiểm chứng lựa
chọn này bằng thí nghiệm.

\subsection{Các mô hình được so sánh}

Bảy phương pháp thuộc bốn nhóm được đặt cạnh nhau (Bảng~\ref{tab:phuongphap}).
Hai mốc so sánh không cần huấn luyện và cho biết mức sai số ``miễn phí''; mọi
phương pháp học phải vượt qua mốc này mới có giá trị.

\begin{table}[H]
\centering
\caption{Các phương pháp được huấn luyện và so sánh.}
\label{tab:phuongphap}
\begin{tabularx}{\textwidth}{>{\bfseries\color{daunavy}}l >{\color{daumaroon}}l X}
\toprule
\rowcolor{gray100}\bfseries\color{daunavy}Nhóm & \bfseries\color{daumaroon}Phương pháp & \leavevmode\bfseries\color{daunavy}Cấu hình \\
\midrule
Mốc so sánh & Persistence & Giữ nguyên vận tốc 6 giờ gần nhất \\
Mốc so sánh & CLIPER-proxy & Persistence có suy giảm \\
AI cổ điển & k-NN analog & $k = 40$; trọng số quỹ đạo 1,0, vị trí 1,0, mùa 0,5 \\
Học máy & Gradient boosting & HistGradientBoosting, hồi quy độ dịch chuyển \\
Học sâu & LSTM & 2 tầng, 128 nơ-ron ẩn, dropout 0,2 \\
Học sâu & Attention-LSTM & LSTM + attention cộng trên chuỗi đầu vào \\
Học sâu & Transformer & Encoder self-attention + mã hóa vị trí \\
\bottomrule
\end{tabularx}
\end{table}

Trong các phương pháp trên, k-NN analog có vai trò riêng: ngoài việc đưa ra
gợi ý, nó chỉ ra những cơn bão lịch sử làm căn cứ cho gợi ý đó.
Giải thuật~\ref{alg:knn} mô tả cách thực hiện.

\begin{algorithm}[H]
\caption{Gợi ý đường đi bằng k-NN analog}
\label{alg:knn}
\hspace*{\algorithmicindent}\textbf{Đầu vào:} cửa sổ quan trắc $q$; thư viện tình huống lịch sử $\mathcal{L}$; $k$; trọng số $w_g$\\
\hspace*{\algorithmicindent}\textbf{Đầu ra:} quỹ đạo gợi ý $\hat{\mathbf{y}}$ và danh sách bão tương tự
\begin{algorithmic}[1]
\Function{GoiYAnalog}{$q, \mathcal{L}, k$}
  \State $\mathbf{p}_t \gets$ vị trí cuối của $q$
  \State Tính đặc trưng $\mathbf{z}_g(q)$: đoạn quỹ đạo tương đối so với $\mathbf{p}_t$, vị trí, mùa
  \For{mỗi tình huống $c \in \mathcal{L}$}
    \State Tính $D(q,c)$ theo công thức~\eqref{eq:knn-dist}
  \EndFor
  \State $N \gets k$ tình huống có $D$ nhỏ nhất
  \State $\alpha_j \gets \dfrac{1/(D_j+\epsilon)}{\sum_{i \in N} 1/(D_i+\epsilon)}$ với mọi $j \in N$
  \State $\hat{\mathbf{y}} \gets \mathbf{p}_t + \sum_{j \in N} \alpha_j \, \Delta\mathbf{y}_j$
  \State \Return $\hat{\mathbf{y}}$, tên và năm của 5 cơn bão khác nhau gần $q$ nhất
\EndFunction
\end{algorithmic}
\end{algorithm}

\subsection{Cấu hình huấn luyện mạng nơ-ron}

\begin{itemize}[leftmargin=1.6em, itemsep=2pt]
  \item Hàm mất mát Smooth~L1 trên độ dịch chuyển (đơn vị độ).
  \item Adam, tốc độ học $10^{-3}$, weight decay $10^{-5}$, lịch cosine,
        cắt gradient ở chuẩn 1,0.
  \item Tối đa 100 epoch, dừng sớm khi validation không cải thiện sau 10 epoch,
        giữ trọng số của epoch tốt nhất trên validation.
  \item Mỗi mạng chạy 3 seed (0, 1, 2), báo cáo trung bình $\pm$ độ lệch chuẩn.
\end{itemize}

\section{Thiết kế phần phục vụ}

\subsection{Gói mô hình}

Kết quả huấn luyện được xuất thành một thư mục gồm trọng số, tham số của bộ
chuẩn hóa, bản kê chia tập và một tệp \code{manifest} ghi phiên bản, cấu hình
kiến trúc, danh sách đặc trưng, các mốc dự báo, bán kính nón và mã băm SHA-256
của từng tệp. Khi nạp, phần phục vụ kiểm tra lần lượt phiên bản lược đồ, mã
băm từng tệp, danh sách đặc trưng và kiểu mô hình thuộc danh sách cho phép;
sai một bước là từ chối nạp. Trọng số chỉ được đọc ở chế độ dữ liệu thuần,
không thực thi mã đi kèm tệp.

\subsection{Giao diện API}

Phần phục vụ là một ứng dụng FastAPI. Mỗi phản hồi dự báo mang theo phiên bản
mô hình, phiên bản dữ liệu, nguồn quan trắc, phương pháp tính bất định, các
cảnh báo (ví dụ quan trắc cũ hơn 24 giờ) và dòng tuyên bố miễn trừ. Các lỗi
trả về theo một khuôn dạng thống nhất gồm mã lỗi, thông điệp và mã yêu cầu để
tra nhật ký. Những đường dẫn ghi dữ liệu hoặc chạy mô hình yêu cầu khóa API
hoặc token, và các đường dẫn dữ liệu đều bị giới hạn tần suất theo địa chỉ
gọi.

\subsection{Nguyên tắc an toàn khi phục vụ}

\begin{itemize}[leftmargin=1.6em, itemsep=2pt]
  \item Ở chế độ vận hành, hệ thống từ chối khởi động nếu thiếu thông tin xác
        thực hoặc vẫn cho phép phương pháp dự phòng, thay vì lặng lẽ chạy
        không có mô hình đã duyệt.
  \item Đầu vào bị từ chối nếu quan trắc không theo thứ tự thời gian, có
        thời điểm ở tương lai hoặc cho tốc độ di chuyển phi thực tế.
  \item Bộ nhớ đệm chỉ là phần tăng tốc: khi nó hỏng, dự báo vẫn được tính
        lại từ đầu và kết quả không đổi.
\end{itemize}

\section{Quy trình chọn mô hình và nón bất định}

Mọi lựa chọn, gồm kiến trúc, đặc trưng, siêu tham số và thành viên của mô
hình kết hợp, chỉ dựa trên tập validation 2017--2020. Tập kiểm thử 2021--2024
được chạy \textbf{một lần} sau khi đã chốt cấu hình. Bán kính nón bất định
tại mỗi mốc lấy bằng phân vị 80\% của sai số trên validation.

\begin{luuy}
Gradient boosting có dòng dẫn đường đạt điểm test tốt nhất trong một thí
nghiệm, nhưng kém hơn trên validation nên không được chọn. Chọn mô hình theo
điểm test sẽ làm con số báo cáo lạc quan hơn thực tế.
\end{luuy}

% ---------------------------------------------------------------------
%  CHƯƠNG 4 — CÀI ĐẶT VÀ THỬ NGHIỆM
% ---------------------------------------------------------------------
\chapter{Cài đặt và thử nghiệm}
\label{ch:caidat}

\begin{muctieu}
Chương này mô tả môi trường triển khai, các thành phần chính kèm mã nguồn
minh hoạ, quy trình huấn luyện, giao diện demo và cách kiểm thử hệ thống.
\end{muctieu}

\section{Môi trường triển khai}

Hệ thống được phát triển và huấn luyện trên máy tính cá nhân chạy Windows 11,
Python 3.11, PyTorch bản CPU cùng các thư viện ở Bảng~\ref{tab:congcu}. Mạng
nơ-ron có kích thước nhỏ (đầu vào 4 bước, 2 tầng LSTM 128 nơ-ron) nên huấn
luyện được trên CPU, không cần GPU. Phụ thuộc được khai báo trong
\code{pyproject.toml}, chia nhóm \code{train}, \code{api}, \code{ingestion},
\code{dev}.

Mã nguồn được quản lý bằng \code{git} và đặt tại kho:
\begin{center}
\url{https://github.com/VietManh05/typhoon-track-vietnam}
\end{center}

\textbf{Sử dụng công cụ trí tuệ nhân tạo hỗ trợ:} nhóm dùng trợ lý lập trình
Claude (Anthropic) để hỗ trợ viết một phần mã nguồn và soạn thảo báo cáo LaTeX.
Mọi số liệu trong báo cáo được sinh ra bằng cách chạy mã nguồn của dự án, và
nhóm đã kiểm tra lại toàn bộ nội dung.

\section{Các thành phần chính}

\subsection{Mô hình LSTM dự báo độ dịch chuyển}

Mã nguồn~\ref{lst:lstm} là hàm lan truyền xuôi của \code{LSTMTrackForecaster}
(tệp \code{src/typhoon\_vn/models/lstm.py}). Đầu ra hồi quy được cộng với vị trí
quan trắc cuối (đã đổi về độ thật), hiện thực hoá thiết kế dự báo độ dịch chuyển.

\begin{lstlisting}[language=Python, caption={Lan truyền xuôi của mô hình LSTM (rút gọn, bỏ các đầu ra gió và khí áp).}, label={lst:lstm}]
def forward(self, x, mask=None):
    # x: (batch, time, features)
    context = self.encode(x, mask)  # last hidden state
    H = self.n_horizons
    reg = self.reg_head(context).view(-1, H, 2)
    if self.residual:
        _, lengths = _validate_sequence(x, mask)
        # last observed lat/lon, scaled
        last = _last_valid(x[..., :2], lengths)
        # back to physical degrees
        anchor = (last * self.position_scale
                  + self.position_mean)
        # position + predicted displacement
        reg = anchor.unsqueeze(1) + reg
    cls = self.cls_head(context).view(-1, H, self.n_classes)
    return {"reg": reg, "cls": cls}
\end{lstlisting}

\subsection{k-NN analog và lời giải thích}

Mã nguồn~\ref{lst:knn} cài đặt Giải thuật~\ref{alg:knn}
(tệp \code{src/typhoon\_vn/models/classical.py}): lấy $k$ láng giềng, đánh
trọng số nghịch đảo khoảng cách, cộng độ dịch chuyển trung bình vào vị trí cuối.
Hàm \code{explain} trả về các cơn bão lịch sử khác nhau gần nhất để hiển thị.

\begin{lstlisting}[language=Python, caption={Dự báo của k-NN analog.}, label={lst:knn}]
def predict(self, x_raw):
    """Forecast positions (N, H, 2) in degrees."""
    distance, index = self._neighbours(x_raw)
    # inverse-distance weights, normalised per forecast
    weight = 1.0 / (distance + 1e-6)
    weight /= weight.sum(axis=1, keepdims=True)
    targets = self._targets[index]
    displacement = np.einsum("nk,nkhc->nhc", weight, targets)
    return anchors(x_raw)[:, None, :] + displacement
\end{lstlisting}

\subsection{Vòng huấn luyện}

Lớp \code{Trainer} (tệp \code{src/typhoon\_vn/training/trainer.py}) thực hiện
vòng lặp epoch, cắt gradient, lịch học cosine, dừng sớm theo validation, lưu
checkpoint và cho phép huấn luyện tiếp từ checkpoint (khôi phục cả trạng thái
bộ sinh số ngẫu nhiên để tái lập kết quả).

\subsection{Đóng gói mô hình}

Sau huấn luyện, mô hình được đóng gói thành thư mục gồm \code{weights.pt},
\code{scaler.json}, \code{split-manifest.json} và \code{manifest.json}. Tệp
manifest lưu phiên bản, mã băm SHA-256 của tập dữ liệu, cấu hình mô hình, mã
băm danh sách đặc trưng, bán kính nón và checksum từng tệp. Khi nạp, hệ thống
kiểm tra toàn bộ checksum, số đặc trưng, số mốc dự báo và nạp trọng số với
\code{weights\_only=True}; gói bị sửa hoặc sai định dạng sẽ bị từ chối.

\subsection{Chạy các thí nghiệm}

Mã nguồn~\ref{lst:run} là các lệnh tái lập toàn bộ kết quả ở
Chương~\ref{ch:ketqua}.

\begin{lstlisting}[language=bash, caption={Các lệnh tái lập thí nghiệm.}, label={lst:run}]
# Benchmark on IBTrACS and uncertainty-cone calibration
python scripts/run_ibtracs_benchmark.py --data-root data \
    --output-dir runs/ibtracs-benchmark
# Compare all methods with three seeds
python scripts/run_model_comparison.py --data-root data \
    --output-dir runs/model-comparison --seeds 0 1 2
# Error-by-horizon figure and case studies
python scripts/plot_error_analysis.py --run runs/model-comparison \
    --data-root data
# Build the analog library used by the API
python scripts/build_analog_artifact.py --data-root data \
    --output runs/analog/analog-wp.npz
\end{lstlisting}

\section{Giao diện và demo}

Mô hình được phục vụ qua FastAPI với các điểm cuối chính \code{/forecast},
\code{/typhoons/active} và \code{/typhoons/\{id\}/track} (quỹ đạo dự báo kèm
nón bất định). Giao diện bản đồ web hiển thị quỹ đạo thực, quỹ đạo dự báo của
LSTM, nón bất định, quỹ đạo gợi ý của k-NN analog (đường màu tím) và khung
``Gợi ý theo bão tương tự'' liệt kê 5 cơn bão lịch sử giống nhất. Thư viện
analog chỉ gồm bão 1980--2020, không chứa bão kiểm thử.

Bản demo công khai tại \url{https://baoviet2026.pages.dev} phát lại 20 cơn bão
thật giai đoạn 2021--2024 ảnh hưởng Việt Nam, tại thời điểm vị trí +72 giờ gần
bờ biển nhất, kèm quỹ đạo thực tế và sai số. Hình~\ref{fig:yagi} là ví dụ cho
bão YAGI 2024.

\begin{figure}[H]
\centering
\includegraphics[width=0.85\textwidth]{case_WP122024.png}
\caption{Quỹ đạo dự báo của các phương pháp cho bão YAGI (phát lúc 06Z 04/09/2024).}
\label{fig:yagi}
\end{figure}

\begin{luuy}
Khi phát lại một cơn bão cũ với thời gian dời về hiện tại, đặc trưng mùa của
k-NN thay đổi (ví dụ YAGI tháng 9 thành tháng 11) và cho ra bão tương tự khác
hẳn. Vì vậy bản demo giữ đúng ngày lịch sử của cơn bão.
\end{luuy}

\section{Kiểm thử}

Dự án có bộ kiểm thử tự động bằng \code{pytest} (thư mục \code{tests/}), chia
thành ba nhóm:
\begin{itemize}[leftmargin=1.6em, itemsep=2pt]
  \item \textbf{Dữ liệu:} parser các nguồn, chuẩn hóa mã bão, gộp đa nguồn, ghi
        Parquet; dữ liệu biên như tọa độ ngoài phạm vi, thời điểm thiếu múi giờ,
        giá trị gió không hữu hạn đều bị từ chối.
  \item \textbf{Học máy:} bộ chuẩn hóa chỉ fit trên train, chia tập không trùng
        cơn bão, kích thước đầu ra mô hình, huấn luyện tiếp từ checkpoint, gói mô
        hình bị sửa checksum phải bị từ chối.
  \item \textbf{API:} kiểm tra đầu vào (ít nhất 2 quan trắc, khoảng cách thời
        gian hợp lệ, tốc độ suy ra không quá 150~km/h) và các điểm cuối.
\end{itemize}

Một lỗi phát hiện trong quá trình làm là timestamp trùng hoặc lùi trong dữ
liệu gốc sinh ra tốc độ và gia tốc giả rất lớn; lỗi được khắc phục bằng cách
không gán khoảng thời gian 6 giờ giả định cho các bước này.

% ---------------------------------------------------------------------
%  CHƯƠNG 5 — KẾT QUẢ VÀ ĐÁNH GIÁ
% ---------------------------------------------------------------------
\chapter{Kết quả và đánh giá}
\label{ch:ketqua}

\begin{muctieu}
Chương này đối chiếu kết quả với các mục tiêu đã đặt ra, sau đó trình bày
chi tiết sai số của từng phương pháp, kiểm định thống kê, độ phủ của nón bất
định, hai nhóm thí nghiệm bổ sung, kết quả kiểm thử phần mềm, phần bàn luận
và những hạn chế còn tồn tại.
\end{muctieu}

\section{Kết quả đạt được}

Bảng~\ref{tab:muctieu} đối chiếu từng mục tiêu đã nêu ở Chương~\ref{ch:modau}.
Bốn mục tiêu đầu đạt đầy đủ; mục tiêu triển khai đạt một phần vì mô hình tốt
nhất trong thí nghiệm chưa phải là mô hình đang phục vụ. Các mục sau trình bày
bằng chứng cho từng dòng của bảng.

\begin{table}[H]
\centering
\caption{Đối chiếu mục tiêu đồ án.}
\label{tab:muctieu}
\begin{tabularx}{\textwidth}{>{\bfseries\color{daunavy}}l X >{\leavevmode\color{daumaroon}\raggedright\arraybackslash}p{3.4cm}}
\toprule
\rowcolor{gray100}\bfseries\color{daunavy}Mục tiêu & \leavevmode\bfseries\color{daunavy}Kết quả & \leavevmode\bfseries\color{daumaroon}Mức độ \\
\midrule
MT1 & Quy trình dữ liệu 24.024 quan trắc, chia theo năm, scaler fit trên train & Đạt \\
MT2 & 7 phương pháp, 3 seed, bootstrap theo cơn bão & Đạt \\
MT3 & 72 giờ: 676~km (persistence) xuống 462~km (k-NN), giảm 32\% & Đạt \\
MT4 & Nón 80\% phủ 82--86\% vị trí thực trên test & Đạt \\
MT5 & API, bản đồ web, demo công khai kèm bão tương tự; mô hình đang phục vụ chưa dùng dòng dẫn đường & Đạt một phần \\
\bottomrule
\end{tabularx}
\end{table}

\section{Độ đo đánh giá}

Độ đo chính là sai số Haversine trung bình tại từng mốc dự báo, tính bằng
kilômét theo công thức~\eqref{eq:haversine}. Đây là khoảng cách trên mặt cầu
giữa vị trí gợi ý và vị trí thực, nên càng nhỏ càng tốt. Các mạng nơ-ron chạy
3 seed và được báo cáo dưới dạng trung bình $\pm$ độ lệch chuẩn; persistence,
CLIPER-proxy và k-NN là tất định nên chỉ có một giá trị.

Kết quả được tính trên hai tập. Bảng~\ref{tab:ketqua-all} dùng toàn bộ 84 cơn
bão kiểm thử (714 cửa sổ). Bảng~\ref{tab:ketqua-vn} chỉ dùng 27 cơn bão ảnh
hưởng Việt Nam (278 cửa sổ), là nhóm mà đề tài quan tâm nhất.

\begin{table}[H]
\centering
\caption{Sai số trung bình (km) trên toàn bộ tập kiểm thử 2021--2024.}
\label{tab:ketqua-all}
\begin{tabular*}{\textwidth}{@{\extracolsep{\fill}}>{\bfseries\color{daunavy}}l rrrrr}
\toprule
\rowcolor{gray100}\bfseries\color{daunavy}Phương pháp & \bfseries\color{daumaroon}6h & \bfseries\color{daumaroon}12h & \bfseries\color{daumaroon}24h & \bfseries\color{daumaroon}48h & \bfseries\color{daumaroon}72h \\
\midrule
Persistence       & 31 & 68 & 155 & 390 & 676 \\
CLIPER-proxy      & 30 & 65 & 151 & 388 & 677 \\
k-NN analog       & 32 & 64 & 133 & 290 & \textbf{462} \\
Gradient boosting & \textbf{29} & \textbf{60} & \textbf{129} & 292 & 476 \\
LSTM              & 33 $\pm$ 0 & 65 $\pm$ 1 & 133 $\pm$ 1 & \textbf{289} $\pm$ 1 & 464 $\pm$ 5 \\
Attention-LSTM    & 35 $\pm$ 3 & 68 $\pm$ 4 & 138 $\pm$ 5 & 297 $\pm$ 10 & 473 $\pm$ 9 \\
Transformer       & 33 $\pm$ 1 & 64 $\pm$ 1 & 133 $\pm$ 3 & 293 $\pm$ 3 & 468 $\pm$ 5 \\
\bottomrule
\end{tabular*}
\end{table}

\begin{table}[H]
\centering
\caption{Sai số trung bình (km) trên 27 cơn bão ảnh hưởng Việt Nam.}
\label{tab:ketqua-vn}
\begin{tabular*}{\textwidth}{@{\extracolsep{\fill}}>{\bfseries\color{daunavy}}l rrrrr}
\toprule
\rowcolor{gray100}\bfseries\color{daunavy}Phương pháp & \bfseries\color{daumaroon}6h & \bfseries\color{daumaroon}12h & \bfseries\color{daumaroon}24h & \bfseries\color{daumaroon}48h & \bfseries\color{daumaroon}72h \\
\midrule
Persistence       & 33 & 70 & 159 & 388 & 642 \\
CLIPER-proxy      & 32 & 68 & 159 & 414 & 710 \\
k-NN analog       & 33 & 65 & 134 & 290 & 476 \\
Gradient boosting & \textbf{30} & 61 & 131 & 295 & 491 \\
LSTM              & 32 $\pm$ 0 & 63 $\pm$ 0 & 133 $\pm$ 1 & 292 $\pm$ 6 & 468 $\pm$ 9 \\
Attention-LSTM    & 34 $\pm$ 3 & 66 $\pm$ 4 & 138 $\pm$ 4 & 298 $\pm$ 6 & 471 $\pm$ 3 \\
Transformer       & 32 $\pm$ 1 & \textbf{62} $\pm$ 1 & \textbf{130} $\pm$ 2 & \textbf{289} $\pm$ 3 & \textbf{463} $\pm$ 10 \\
\bottomrule
\end{tabular*}
\end{table}

Hai bảng cho cùng một bức tranh. Ở 6 và 12 giờ, mọi phương pháp chỉ cách
nhau vài kilômét. Từ 24 giờ, nhóm phương pháp học tách khỏi hai mốc so sánh,
và đến 72 giờ khoảng cách lên tới khoảng 200~km. Trong nhóm phương pháp học,
không có phương pháp nào dẫn đầu ở mọi mốc. Hình~\ref{fig:horizon} biểu diễn
xu hướng này: các đường gần như trùng nhau ở 6 giờ và tách dần từ 24 giờ trở
đi.

\begin{figure}[H]
\centering
\includegraphics[width=0.85\textwidth]{error_by_horizon.png}
\caption{Sai số trung bình theo thời hạn dự báo của các phương pháp.}
\label{fig:horizon}
\end{figure}

\subsection{Kiểm định thống kê}

Chênh lệch vài chục kilômét giữa hai phương pháp có thể chỉ do may rủi của
tập kiểm thử. Để kiểm tra, Bảng~\ref{tab:bootstrap} so sánh tại 72 giờ bằng
bootstrap theo cơn bão: 50 cơn bão kiểm thử có dự báo 72 giờ được lấy mẫu lại
5.000 lần. Đơn vị lấy mẫu là cả cơn bão chứ không phải từng cửa sổ, vì các
cửa sổ của cùng một cơn bão không độc lập với nhau. Giá trị âm nghĩa là tốt
hơn mốc tham chiếu.

\begin{table}[H]
\centering
\caption{Chênh lệch sai số 72 giờ (km) và khoảng tin cậy 95\%.}
\label{tab:bootstrap}
\begin{tabular*}{\textwidth}{@{\extracolsep{\fill}}>{\bfseries\color{daunavy}}l cc}
\toprule
\rowcolor{gray100}\bfseries\color{daunavy}Phương pháp & \bfseries\color{daumaroon}So với persistence & \bfseries\color{daumaroon}So với LSTM \\
\midrule
CLIPER-proxy      & $-15$ \ [$-101$; $+65$]  & $+188$ \ [$+113$; $+256$] \\
k-NN analog       & $-201$ \ [$-268$; $-131$] & $+3$ \ [$-29$; $+34$] \\
Gradient boosting & $-191$ \ [$-252$; $-126$] & $+12$ \ [$-20$; $+48$] \\
LSTM              & $-203$ \ [$-262$; $-145$] & --- \\
Attention-LSTM    & $-173$ \ [$-229$; $-116$] & $+30$ \ [$+12$; $+50$] \\
Transformer       & $-200$ \ [$-257$; $-142$] & $+3$ \ [$-16$; $+22$] \\
\bottomrule
\end{tabular*}
\end{table}

Bảng cho ba kết luận. Mọi phương pháp học đều tốt hơn persistence 170--200~km
với khoảng tin cậy không chứa 0, nên mục tiêu MT3 đạt một cách chắc chắn.
Giữa LSTM, Transformer, k-NN và gradient boosting không có khác biệt có ý
nghĩa thống kê. Riêng Attention-LSTM kém hơn LSTM có ý nghĩa, và CLIPER-proxy
không khác persistence.

\subsection{Nón bất định}

Bán kính nón, lấy bằng phân vị 80\% của sai số trên validation, là 52, 97,
210, 455 và 714~km tại 6, 12, 24, 48 và 72 giờ. Trên tập kiểm thử, tỉ lệ vị
trí thực nằm trong nón là 82--86\% ở các mốc. Con số này sát mức danh nghĩa
80\% và hơi cao hơn, tức nón được hiệu chỉnh tốt và nghiêng về phía thận
trọng. Cần lưu ý nón ở 72 giờ rộng hơn 700~km: nó trung thực về mức độ bất
định, nhưng cũng cho thấy gợi ý ở hạn dài chỉ nên dùng để định hướng.

\subsection{Thí nghiệm loại bỏ: dự báo độ dịch chuyển}
\label{sec:ablation}

Chương~\ref{ch:thietke} đã lập luận rằng nên cho mạng học độ dịch chuyển thay
vì tọa độ tuyệt đối. Thí nghiệm này kiểm chứng lập luận đó: giữ nguyên kiến
trúc LSTM, dữ liệu và seed 0, chỉ đổi mục tiêu học (Bảng~\ref{tab:ablation}).

\begin{table}[H]
\centering
\caption{Ảnh hưởng của cách biểu diễn đầu ra lên sai số LSTM (km).}
\label{tab:ablation}
\begin{tabular*}{\textwidth}{@{\extracolsep{\fill}}>{\bfseries\color{daunavy}}l rrrrr}
\toprule
\rowcolor{gray100}\bfseries\color{daunavy}Mục tiêu học & \bfseries\color{daumaroon}6h & \bfseries\color{daumaroon}12h & \bfseries\color{daumaroon}24h & \bfseries\color{daumaroon}48h & \bfseries\color{daumaroon}72h \\
\midrule
Tọa độ tuyệt đối         & 115 & 121 & 164 & 305 & 473 \\
Độ dịch chuyển (đang dùng) & \textbf{33} & \textbf{65} & \textbf{133} & \textbf{289} & \textbf{464} \\
\bottomrule
\end{tabular*}
\end{table}

Khác biệt rõ nhất ở hạn ngắn. Học tọa độ tuyệt đối làm sai số 6 giờ tăng
khoảng 3,5 lần, kém cả persistence (31~km), và cần nhiều epoch hơn mới hội tụ
(dừng ở epoch 42 thay vì 13--17). Ở 72 giờ hai cách gần bằng nhau, vì khi đó
sai số do hướng đi tương lai quyết định chứ không do vị trí xuất phát. Kết
quả này xác nhận lựa chọn thiết kế: phần mạng cần học là chuyển động, không
phải vị trí hiện tại.

\subsection{Tối ưu huấn luyện}
\label{sec:steering}

Sau khi các phương pháp học cho kết quả ngang nhau, câu hỏi tiếp theo là còn
có thể giảm sai số bằng cách nào. Hai hướng được thử: thay đổi cách huấn
luyện trên cùng thông tin đầu vào, và bổ sung thông tin đầu vào mới. Các thí
nghiệm này đánh giá trên cùng các cửa sổ của bão vùng Việt Nam có đủ 48 giờ
lịch sử (435 cửa sổ validation, 526 cửa sổ test). Mọi lựa chọn dựa trên
validation; test chạy một lần.

\textbf{(a) Chỉ dùng thông tin quỹ đạo: đã chạm trần.} LSTM gốc có sai số
trung bình trên validation là 201~km. Không thay đổi nào cải thiện được con
số này: dùng dữ liệu toàn lưu vực (212~km), lịch sử 48 giờ (212~km), hàm mất
mát theo km (203~km), mạng lớn hoặc nhỏ hơn (206--207~km), Transformer
(207~km). Mô hình kết hợp LSTM và Transformer trên test chỉ giảm 8~km tại 72
giờ, khoảng tin cậy [$-16$; $+2$], không có ý nghĩa thống kê.

\textbf{(b) Thêm dòng dẫn đường khí quyển: cải thiện có ý nghĩa.} Hướng thứ
hai bổ sung 11 đặc trưng từ NCEP/NCAR R1 tại đúng thời điểm quan trắc
(Bảng~\ref{tab:dactrung}). Kết quả ở Bảng~\ref{tab:steering}.

\begin{table}[H]
\centering
\caption{Hiệu quả của đặc trưng dòng dẫn đường trên tập kiểm thử (km, 3 seed).}
\label{tab:steering}
\begin{tabular*}{\textwidth}{@{\extracolsep{\fill}}>{\bfseries\color{daunavy}}l rrrrrr}
\toprule
\rowcolor{gray100}\bfseries\color{daunavy}Cấu hình & \bfseries\color{daumaroon}6h & \bfseries\color{daumaroon}12h & \bfseries\color{daumaroon}24h & \bfseries\color{daumaroon}48h & \bfseries\color{daumaroon}72h & \bfseries\color{daumaroon}TB \\
\midrule
LSTM chỉ quỹ đạo (gốc)        & 32 & 64 & 131 & 286 & 469 & 197 \\
LSTM + dòng dẫn đường         & 30 & 58 & 118 & 270 & 465 & 188 \\
Transformer + dòng dẫn đường  & 30 & 57 & 115 & 265 & 450 & 183 \\
Kết hợp LSTM + Transformer    & \textbf{28} & \textbf{54} & \textbf{112} & \textbf{259} & \textbf{443} & \textbf{179} \\
\bottomrule
\end{tabular*}
\end{table}

Mô hình kết hợp (chọn trên validation) giảm 35~km tại 72 giờ so với mô hình
gốc, khoảng tin cậy 95\% [$-62$; $-9$] trên 44 cơn bão, có ý nghĩa thống kê;
tại 24 giờ giảm khoảng 15\%, trung bình các mốc giảm khoảng 9\%.

\section{Kiểm thử phần mềm}

Ngoài đánh giá mô hình, phần mềm được kiểm tra bằng bộ kiểm thử tự động viết
với \code{pytest}. Bộ kiểm thử bao phủ các khâu nạp dữ liệu, hợp đồng đặc
trưng, vòng huấn luyện, đóng gói và nạp gói mô hình, suy luận, API, lưu trữ,
bộ nhớ đệm và tiến trình cập nhật. Ở phiên bản bàn giao, 269 ca kiểm thử chạy
đạt và 1 ca được bỏ qua vì cần một máy chủ PostGIS bên ngoài.

Dữ liệu dùng trong kiểm thử phần mềm là dữ liệu tổng hợp. Các ca này chỉ
chứng minh phần mềm hoạt động đúng hợp đồng, không phải bằng chứng về độ
chính xác dự báo; độ chính xác chỉ được kết luận từ các bảng ở mục trên.

\section{Bàn luận}

\textbf{Học từ dữ liệu giúp rõ ở hạn dài.} Ở 6 giờ persistence gần như tối ưu,
vì trong thời gian ngắn như vậy bão chưa kịp đổi hướng. Từ 24 giờ trở đi, các
phương pháp học giảm sai số khoảng 14\% (24 giờ) đến 30\% (72 giờ). Lý do là
chúng học được các kiểu quỹ đạo điển hình, chẳng hạn đi về phía tây rồi
chuyển hướng lên phía bắc, điều mà phép ngoại suy thẳng không thể biết.

\textbf{Mô hình sâu không vượt trội.} Với khoảng 7.800 cửa sổ huấn luyện và 4
bước đầu vào, k-NN analog và gradient boosting đạt cùng mức với LSTM và
Transformer. Mạng nơ-ron overfit nhanh, epoch tốt nhất chỉ là 2--9. Thí
nghiệm (a) cho thấy đổi kiến trúc hay siêu tham số không giúp gì, còn thí
nghiệm (b) cho thấy thêm thông tin khí quyển thì có. Từ đó rút ra kết luận
chính của đồ án: nút thắt là thông tin đầu vào, không phải sức chứa của mô
hình. Điều này hợp với vật lý, vì bão bị dòng dẫn đường đẩy đi, và dữ liệu
chỉ có quỹ đạo không chứa thông tin này.

\textbf{Sai số theo cường độ.} Tại 72 giờ, áp thấp nhiệt đới là nhóm khó nhất
với persistence (800~km), còn bão rất mạnh cấp 12--13 là nhóm dễ nhất (khoảng
400~km với các phương pháp học). Một giả thuyết cần kiểm chứng thêm là hướng
đi của áp thấp kém ổn định hơn.

\textbf{Nghiên cứu trường hợp.} Sai số trung bình che mất những ca khó. Vì
vậy Bảng~\ref{tab:case} chọn thời điểm phát dự báo sao cho vị trí sau 72 giờ
gần bờ biển Việt Nam nhất, tức trường hợp bão sắp đổ bộ và gợi ý có giá trị
sử dụng cao nhất.

\begin{table}[H]
\centering
\caption{Sai số 72 giờ (km) khi bão sắp đổ bộ Việt Nam.}
\label{tab:case}
\small
\begin{tabular*}{\textwidth}{@{\extracolsep{\fill}}>{\bfseries\color{daunavy}}l rrrrr}
\toprule
\rowcolor{gray100}\bfseries\color{daunavy}Bão (thời điểm phát) & \bfseries\color{daumaroon}k-NN & \bfseries\color{daumaroon}GB & \bfseries\color{daumaroon}LSTM & \bfseries\color{daumaroon}Transf. & \bfseries\color{daumaroon}Persist. \\
\midrule
YAGI 2024 (06Z 04/09)    & 419 & \textbf{407} & 541 & 678 & 701 \\
SAOLA 2023 (12Z 31/08)   & \textbf{365} & 538 & 534 & 514 & 534 \\
YINXING 2024 (06Z 09/11) & 575 & \textbf{269} & 700 & 550 & 958 \\
CONSON 2021 (00Z 10/09)  & 558 & 581 & 625 & 702 & 987 \\
\bottomrule
\end{tabular*}
\end{table}

Với YAGI (Hình~\ref{fig:yagi}), phần lớn mô hình gợi ý bão chuyển hướng lên
phía bắc, kiểu quỹ đạo phổ biến trong dữ liệu, trong khi YAGI đi thẳng về phía
tây vào miền Bắc Việt Nam. Gradient boosting (407~km) và k-NN analog (419~km)
gần đúng nhất; k-NN dựa trên các bão tương tự như WUKONG 2000, DAMREY 2005. Với SAOLA, k-NN tốt nhất (365~km so với
514--579~km của các mạng nơ-ron). Với CONSON, CLIPER-proxy lại tốt nhất
(101~km); đây là trường hợp đơn lẻ, không đủ để kết luận. Các ví dụ này cho thấy
giá trị của k-NN không chỉ ở độ chính xác mà ở chỗ người dùng thấy được gợi ý
dựa trên những cơn bão nào.

\section{Hạn chế}

\begin{luuy}
\begin{itemize}[leftmargin=1.4em, itemsep=2pt]
  \item Mô hình đang phục vụ trên website và API vẫn là LSTM chỉ dùng quỹ đạo;
        mô hình có dòng dẫn đường cần dữ liệu khí quyển tại thời điểm dự báo
        (ví dụ phân tích GFS) mới đưa vào vận hành được.
  \item YAGI 2024 vẫn sai 620--670~km tại 72 giờ với mọi mô hình có dòng dẫn
        đường; lưới 2,5° của NCEP/NCAR và các đặc trưng tóm tắt chưa đủ để nhận
        ra áp cao cận nhiệt mạnh trong ca này.
  \item Nón bất định là đường tròn đẳng hướng, trong khi sai số dọc và ngang
        quỹ đạo thực tế khác nhau.
  \item Kết quả chỉ trên một cách chia theo năm; chưa kiểm định chéo theo mùa.
  \item Dữ liệu bão đang hoạt động từ IBTrACS ACTIVE trễ vài ngày.
  \item Hệ thống chưa được nghiệm thu trong môi trường vận hành thật: chưa có
        bằng chứng về tải, sao lưu và khôi phục trên máy chủ đích.
\end{itemize}
\end{luuy}

% ---------------------------------------------------------------------
%  CHƯƠNG 6 — KẾT LUẬN VÀ HƯỚNG PHÁT TRIỂN
% ---------------------------------------------------------------------
\chapter{Kết luận và hướng phát triển}
\label{ch:ketluan}

\begin{muctieu}
Chương này tổng kết những gì đồ án đã làm được, rút ra bài học và đề xuất
hướng phát triển tiếp theo.
\end{muctieu}

\section{Kết luận}

Đồ án đã xây dựng một hệ thống trí tuệ nhân tạo gợi ý đường đi của cơn bão
ảnh hưởng Việt Nam, đi trọn từ dữ liệu IBTrACS 1980--2024, tạo đặc trưng,
huấn luyện và so sánh bảy phương pháp, đến triển khai API và bản đồ web. Bốn
trên năm mục tiêu đạt đầy đủ (Bảng~\ref{tab:muctieu}): các phương pháp học giảm
sai số 72 giờ khoảng 200~km so với persistence với ý nghĩa thống kê, và nón bất
định phủ 82--86\% vị trí thực. Mục tiêu triển khai mới đạt một phần vì mô hình
tốt nhất, có dòng dẫn đường, chưa được đưa vào phục vụ.

Về mặt huấn luyện, đồ án rút ra ba kết quả chính: dự báo độ dịch chuyển tốt hơn
hẳn dự báo tọa độ tuyệt đối ở hạn ngắn; với dữ liệu chỉ có quỹ đạo, LSTM,
Transformer, k-NN và gradient boosting tương đương nhau; và bổ sung thông tin
dòng dẫn đường khí quyển là thay đổi duy nhất cải thiện có ý nghĩa (giảm 35~km
tại 72 giờ).

\section{Bài học rút ra}

Bài học lớn nhất là \textbf{thông tin đầu vào quan trọng hơn kiến trúc mô hình}.
Chúng em đã thử nhiều cách làm mạng phức tạp hơn mà không cải thiện; chỉ khi
đưa vào dữ liệu phản ánh đúng nguyên nhân vật lý của chuyển động bão thì sai số
mới giảm. Nếu làm lại, chúng em sẽ tìm hiểu về khí tượng trước để chọn đặc trưng, thay vì
bắt đầu từ việc thử kiến trúc.

Bài học thứ hai là về \textbf{kỷ luật đánh giá}: chia tập theo năm, chọn mô hình
chỉ trên validation, chạy nhiều seed và kiểm định bootstrap. Nhiều khác biệt vài
chục km trông có vẻ đáng kể hoá ra không có ý nghĩa thống kê. Chúng em cũng phải
bỏ qua một cấu hình có điểm test tốt nhất vì nó không được chọn trên validation.

Về vai trò của trí tuệ nhân tạo: AI học được các kiểu quỹ đạo phổ biến và giảm
sai số rõ rệt so với ngoại suy đơn giản, nhưng vẫn thất bại ở các cơn bão đi
khác thường như YAGI, đúng những trường hợp nguy hiểm nhất. Vì vậy hệ thống chỉ
nên là công cụ gợi ý bổ sung, có giải thích (bão tương tự) và có vùng bất định,
không thay thế bản tin của cơ quan khí tượng có thẩm quyền.

\section{Hướng phát triển}

\begin{itemize}[leftmargin=1.6em, itemsep=2pt]
  \item Đưa mô hình kết hợp có dòng dẫn đường vào vận hành, lấy đặc trưng từ
        phân tích GFS tại thời điểm dự báo, khắc phục hạn chế mô hình đang phục
        vụ chỉ dùng quỹ đạo.
  \item Dùng tái phân tích ERA5 độ phân giải 0,25° và đặc trưng dạng ảnh (mạng
        CNN trên trường gió) để nhận diện áp cao cận nhiệt, nhắm vào các ca khó
        như YAGI.
  \item Thay nón tròn bằng nón bất đối xứng theo hướng dọc và ngang quỹ đạo,
        hiệu chỉnh riêng theo cường độ bão.
  \item Kiểm định chéo theo nhiều cách chia năm và theo dõi sai số trên các cơn
        bão mới 2025--2026.
\end{itemize}


% ---------------------------------------------------------------------
%  TÀI LIỆU THAM KHẢO VÀ PHỤ LỤC
% ---------------------------------------------------------------------
\backmatter

% \phantomsection tạo mốc neo riêng cho mục lục, nếu thiếu thì hyperref
% báo trùng neo giữa dấu trang và dấu trang cha.
\cleardoublepage
\phantomsection
\addcontentsline{toc}{chapter}{Tài liệu tham khảo}
\bibliographystyle{plain}
\bibliography{refs}

% ---------------------------------------------------------------------
%  PHỤ LỤC
% ---------------------------------------------------------------------
% Trong \backmatter thì \chapter đã tự thêm vào mục lục, không cần
% \addcontentsline nữa, thêm vào là trùng neo dấu trang.
\chapter{Phụ lục}
\markboth{Phụ lục}{}

\section*{A. Hướng dẫn cài đặt và chạy chương trình}
\phantomsection
\addcontentsline{toc}{section}{A. Hướng dẫn cài đặt và chạy chương trình}

Các bước dưới đây cho phép cài đặt, kiểm thử, huấn luyện và chạy hệ thống
từ đầu trên một máy Windows có Python 3.11.

\begin{itemize}[leftmargin=1.6em, itemsep=2pt]
  \item Mã nguồn đầy đủ:\\ \url{https://github.com/VietManh05/typhoon-track-vietnam}
  \item Bản demo trực tuyến: \url{https://baoviet2026.pages.dev}
\end{itemize}

\begin{lstlisting}[language=bash, caption={Các bước cài đặt, huấn luyện và khởi chạy hệ thống.}]
# 1. Clone the repository
git clone https://github.com/VietManh05/typhoon-track-vietnam.git
cd typhoon-track-vietnam

# 2. Create a virtual environment and install dependencies
python -m venv .venv
.venv\Scripts\activate
python -m pip install -e ".[api,train,ingestion,dev]"

# 3. Run the test suite
python -m pytest -q

# 4. Train and compare models (see Chapter 4 for all experiments)
python scripts/run_model_comparison.py --data-root data \
    --output-dir runs/model-comparison --seeds 0 1 2

# 5. Launch the API, then open http://localhost:8000/docs
python -m uvicorn typhoon_vn.api.app:app --reload
\end{lstlisting}

\section*{B. Bảng tự kiểm trước khi nộp}
\phantomsection
\addcontentsline{toc}{section}{B. Bảng tự kiểm trước khi nộp}

Nhóm đã rà soát báo cáo theo danh sách dưới đây trước khi nộp.

\begin{tukiem}
  \item Sản phẩm chạy được trên máy khác, không chỉ trên máy của nhóm.
  \item Đã chạy thử demo trên đúng môi trường sẽ dùng khi bảo vệ.
  \item Mã nguồn nằm trên kho \code{git}, lịch sử commit cho thấy tiến độ
        trải đều chứ không dồn vào tuần cuối.
  \item Có tệp \code{README} và \code{pyproject.toml} để người khác chạy lại.
  \item Mọi hình, bảng và giải thuật đều có chú thích, có nhãn, và được nhắc
        tới trong thân bài.
  \item Mọi mục tiêu ở Chương~\ref{ch:modau} đều được đối chiếu ở
        Chương~\ref{ch:ketluan}.
  \item Mọi con số trong báo cáo trùng khớp với kết quả chạy thật, không có
        số nào là ước lượng hay bịa.
  \item Mọi tài liệu trong \code{refs.bib} đều được trích dẫn ít nhất một
        lần, và mọi nguồn dữ liệu, mã tham khảo đều đã ghi rõ.
  \item Đã khai báo phần nào có dùng công cụ trí tuệ nhân tạo hỗ trợ, và nhóm
        giải thích được toàn bộ hệ thống khi bị hỏi.
  \item Đã xoá hết hộp gợi ý, mọi khối nội dung mẫu, và cả
        chương sổ tra cứu quy cách trình bày.
\end{tukiem}

\section*{C. Nhận xét của giảng viên hướng dẫn}
\phantomsection
\addcontentsline{toc}{section}{C. Nhận xét của giảng viên hướng dẫn}

\khungnhanxet

\vspace{0.8em}
\noindent\begin{minipage}[t]{0.5\textwidth}
\color{daunavy}
{\color{daumaroon}\bfseries Điểm số}

\vspace{0.7em}
\dongchamnhan{Bằng số:}

\vspace{0.9em}
\dongchamnhan{Bằng chữ:}
\end{minipage}\hfill
\okyten{\dauGVHD}{Giảng viên hướng dẫn}


\end{document}
